\documentclass[10pt,a4paper,openany,twoside, prof]{book}

\usepackage{layout_cours}
\setlist{leftmargin=5mm}

\begin{document}

\chapnumber{12} % numéro du chapitre
\titre{Equation de droite} % titre du chapitre
\classe{Seconde} % classe
\entetegal
% \makeatletter
% \addcontentsline{toc}{chapter}{Chapitre \@chapnumber - \@titre} % Prepare a table of contents
% \makeatother

% Debut du texte
% **************
% \paragraph*{Objectif du chapitre :}
% \begin{itemize}
% \item  Reconnaitre une équation de droite.
% \item  Tracer une droite d'équation connue et déterminer l'appartenance d'un point à cette droite.
% \item  Déterminer le coefficient directeur, l'ordonnée à l'origine ainsi que l'équation d'une droite à partir de sa représentation graphique.
% \end{itemize}

\section{Equation cartésienne d'une droite}

\begin{defi}{Droite}{}
\begin{minipage}{.5\textwidth}
	La \textbf{droite} du plan $\R^2$, définie par les points $A$ et $B$, est l'ensemble des points $M$ tel que $A$, $B$ et $M$ sont alignés.
\end{minipage}
\begin{minipage}{.5\textwidth}
   \begin{tikzpicture}[line cap=round,line join=round,>=triangle 45,x=0.65cm,y=0.65cm]
      \draw [xstep=0.65cm,ystep=0.65cm] (0,1) grid (10,5);
      \draw [line width=1.2pt](0,2) -- (10,4);
      \draw (9,4) node[above] {\large{$\mathcal{D}$}};

      \draw [color=blue] (4, 2.8) node[above] {\large{$A$}};
      \fill[color=blue] (4, 2.8) circle (2pt);

      \draw [color=blue] (8, 3.6) node[above] {\large{$B$}};
      \fill[color=blue] (8, 3.6) circle (2pt);

      \draw [color=red] (1, 2.2) node[above] {\large{$M$}};
      \fill[color=red] (1, 2.2) circle (2pt);

      \draw [color=red] (2, 2.4) node[above] {\large{$M$}};
      \fill[color=red] (2, 2.4) circle (2pt);

   \end{tikzpicture}
\end{minipage}
\end{defi}

\begin{defi}{Equation cartésienne de droite}{}
   Soit trois coefficients $a,b,c\in\R$ avec $(a;b)\ne (0;0)$ (c'est-à-dire non nuls simultanément). On appelle \textbf{équation cartésienne de droite} l'équation suivante d'inconnue $(x;y)$ :

   $$ax+by+c=0$$
\end{defi}

\begin{prop}{}{}
	Soit une droite $\mathcal{D}$ du plan $\R^2$. Il existe trois coefficients $a,b,c\in\R$ avec $(a;b)\ne (0;0)$ (c'est-à-dire non simultanément nuls) tels que :

   $$M\coordpp{x}{y}\in\mathcal{D} \iff ax+by+c=0$$

   Autrement dit, $M$ appartient à $\mathcal{D}$, si et seulement si, ses coordonnées $\coordpp{x}{y}$ sont solutions de l'équation $ax+by+c=0$.
\end{prop}

\begin{deme}{}{}
   \begin{reponse}
   Soit deux points \textbf{distincts} $A\coordpp{x_A}{y_a}$ et $B\coordpp{x_B}{y_B}$ de la droite $\mathcal{D}$.

   \begin{itemize}
      \item 
   Procédons par équivalence :
   \begin{align*}
      M\in\mathcal{D} 
      &\iff (AM)\; // \; (AB)  
      &\iff \v{AM} \text{ et } \v{AB} \text{ sont colinéaire} &\\
      &\iff \det{\v{AM}}{\v{AB}} = 0 
      &\iff \Det{x-x_A & x_B-x_A \\ y-y_A & y_B-y_A} = 0 \quad&\\
      &\iff (x-x_A)(y_B-y_A) - (x_B-x_A)(y-y_A) = 0
   \end{align*}

   En posant $a=(y_B-y_A)$, $b=-(x_B-x_A)$ et $c=-x_A(y_B-y_A)+y_A(x_B-x_A)$, 
   
   la dernière ligne se réécrit:

   $$ax+by+c=0$$
   
   \item Vérifions que $(a;b)\ne(0;0)$.
   
   Si $(a;b)=(0;0)$ alors $y_B-y_A=0$ et $x_B-x_A=0$. Ce qui est impossible car les points $A$ et $B$ sont distincts.
\end{itemize}
   \end{reponse}
\end{deme}

\begin{ex}{}{}
   Le point $M\coordpp{5}{-2}$ appartient-il à la droite $\mathcal{D}$ d'équation cartésienne $4x+7y-6=0$ \;?

   \begin{reponse}
      On calcule :

         $4\times 5 + 7 \times (-2) - 6 = 20-14-6=0$

      Donc $M\in\mathcal{D}$
   \end{reponse}
\end{ex}

\begin{rem}{}{}
   L'équation cartésienne d'une droite n'est pas unique ! Trouver une autre équation cartésienne de la droite $\mathcal{D}$ de l'exemple ci-dessus : 
   $$\troum{8x+14y-12=0}$$
   \medskip
\end{rem}


\vspace{-5mm}
\section{Vecteur directeur}

\begin{defi}{Vecteur directeur}{}
   \begin{minipage}{.5\textwidth}
      
      Soient $A$ et $B$ deux point d'une droite $\mathcal{D}$ du plan $\R^2$. 
   
   \noindent Tout vecteur $\vecu$ colinéaire à $\v{AB}$ est appelé \textbf{vecteur directeur} de la droite $\mathcal{D}$.
\end{minipage}
\begin{minipage}{.5\textwidth}
   \begin{tikzpicture}[line cap=round,line join=round,>=triangle 45,x=0.65cm,y=0.65cm]
      \draw [xstep=0.65cm,ystep=0.65cm] (0,1) grid (10,5);
      \draw [line width=1.2pt](0,2) -- (10,4);
      \draw (9,4) node[above] {\large{$\mathcal{D}$}};
      \draw [->,line width=1.2pt, red] (2,1.5) -- (5,2.1);
      \draw[red] (5,2.1) node[right] {\large{$\v{u}$}};

      \draw [color=blue] (4, 2.8) node[above] {\large{$A$}};
      \fill[color=blue] (4, 2.8) circle (2pt);

      \draw [color=blue] (8, 3.6) node[above] {\large{$B$}};
      \fill[color=blue] (8, 3.6) circle (2pt);

   \end{tikzpicture}
\end{minipage}
\end{defi}
   

\begin{prop}{Vecteur directeur à partir de l'équation cartésienne}{vecteur_directeur_eq_cartesienne}
Soit trois coefficients $a$, $b$, $c \in \R$, avec $a$ ou $b$ non nul. Soit la droite $\mathcal{D}$, du plan $\R^2$, d'équation cartésienne $ax+by+c=0$.

Le vecteur $\vecu \coordvp{-b}{a}$ est un vecteur directeur de la droite $\mathcal{D}$.
\end{prop}

\begin{dem}{}{} 
\begin{reponse}
   Il faut montrer que $\vecu$ et $\v{AB}$ sont colinéaires. Il suffit de montrer que $\det{\vecu}{\v{AB}}=0$.

   $$\det{\vecu}{\v{AB}} = \Det{-b & x_B-x_A \\ a & y_B-y_A} = -b(y_B-y_A)- a(x_B-x_A) = - \left[a(x_B-x_A)+b(y_B-y_a)\right]$$

   Puisque $A\in\mathcal{D},\quad ax_A+by_A+c=0$. Et de même, $ax_B+by_B+c=0$.

   Par soustraction des égalités, $a(x_B-x_A)+b(y_B-y_a)=0$
   
   \medskip
   Donc $\det{\vecu}{\v{AB}}=0$ et $\vecu$ est bien un vecteur directeur.
\end{reponse}
\end{dem}


\begin{ex}{}{}
   Trouver un vecteur directeur de la droite d'équation cartésienne $-y-3x=7$.
   
   \begin{reponse}
      $\vecu \coordvp{1}{-3}$ convient. Tout comme $\vecv \coordvp{3}{-9}$ et $\vecw \coordvp{-1}{3}$
   \end{reponse}
   
   \end{ex}

   
\begin{methode}{Trouver l'équation cartésienne...}{}
   \textbf{A partir du vecteur directeur $\vecu\coordvp{-b}{a}$ et du point $A\coordpp{x_A}{y_A}$}

   \medskip
   Nous prendrons l'exemple de $\vecu\coordvp{5}{2}$ et $A\coordpp{-1}{2}$
   \begin{enumerate}[label=(\roman*)]
      \item On commence par déterminer les coefficients $a$ et $b$ de l'équation cartésienne à partir du vecteur directeur. 
      
      Dans notre exemple, $a=\troum{2}$ et $b=\troum{-5}$. Donc l'équation cartésienne est de la forme 
      
      $$\troum{2}x+\troum{-5}y+c=0 \quad\text{avec }c\in\R$$

      \item On détermine le coefficient $c$ avec le point $A$. Puisque $A\in\mathcal{D}$, les coordonnées $\coordpp{x_A}{y_A}$ vérifient l'équation.
      
      Dans notre exemple, $2\times\troum{-1}-5\times\troum{2}+c=0$. Donc $c=\troum{12}$.

   \end{enumerate}
\end{methode}
\begin{methode}{Trouver l'équation cartésienne...}{}
   \textbf{A partir du vecteur $\v{AB}$}

   \medskip
   Nous prendrons l'exemple de  $A(2;-1)$ et $B(-1;5)$. Donc $\v{AB}\coordvp{x_B-x_A}{y_B-y_A}=\coordvp{-3}{6}$.

   Soit $M(x;y)$ un point de la droite $\mathcal{D}$, alors les vecteurs $\v{AM}$ et $\v{AB}$ sont colinéaires :
   
   \medskip
   \hspace{5mm}$\v{AM}\coordvp{x_M-x_A}{y_M-y_A}=\coordvp{x-2}{y+1}$
   
   \medskip
   \noindent $\v{AM}$ et $\v{AB}$ colinéaires $\Longleftrightarrow \det{\v{AM}}{\v{AB}}=\troum{0}$
   
   \noindent \phantom{$\v{AM}$ et $\v{AB}$ colinéaires} $\troum{\Longleftrightarrow (-3)\times(y+1)-6\times(x-2)=0}$

   \noindent \phantom{$\v{AM}$ et $\v{AB}$ colinéaires} $\troum{\Longleftrightarrow -3y-3-6x+12=0}$
   
   \noindent \phantom{$\v{AM}$ et $\v{AB}$ colinéaires} $\troum{\Longleftrightarrow -6x-3y+9=0}$
   
   \medskip
   La droite $\mathcal{D}$ a pour équation cartésienne : $2x+y-3=0$
   \end{methode}
\begin{prop}{Droite horizontale, droite verticale}{droite_verticale}
   \begin{itemize}
      \item 
      La droite $\mathcal{D}$ est parallèle à l'axe des abscisses (horizontale), si et seulement s'il existe $d\in\R$, tel que la droite $\mathcal{D}$ admette $y=d$ comme équation cartésienne. Autrement dit, le coefficient $a$ est nul.
      \item 
      La droite $\mathcal{D}$ est parallèle à l'axe des ordonnées (verticale), si et seulement s'il existe $d\in\R$, tel que la droite $\mathcal{D}$ admette $x=d$ comme équation cartésienne. Autrement dit, le coefficient $b$ est nul.
   \end{itemize}
   \begin{center}
   \begin{tikzpicture}[line cap=round,line join=round,>=triangle 45,x=0.65cm,y=0.65cm]
      \begin{axis}[
         x=1.0cm,y=1.0cm,
         axis lines=middle,ymajorgrids=true,
         xmajorgrids=true,
         xmin=-0.5,xmax=7.5,
         ymin=-0.5,ymax=3.5,
         xtick={-0,1,...,7.0},ytick={-0,1,...,3},
         ]
         \draw [line width=1.2pt, blue] (axis cs:1.5,-0.5) -- (axis cs:1.5,3.5);
         \draw[blue] (axis cs:1.5,3) node[left] {\large{$\mathcal{D}$}};
         \draw[blue] (axis cs:1.5,3) node[right] {\large{$x=1.5$}};
         \draw[line width=1.2pt, red] (axis cs:-0.5,2) -- (axis cs:7.5,2);
         \draw[red] (axis cs:6,2) node[below] {\large{$\mathcal{D'}$}};
         \draw[red] (axis cs:6,2) node[above] {\large{$y=2$}};
      \end{axis}
   \end{tikzpicture}
   \end{center}
\end{prop}

\vspace{-5mm}
\section{Équation réduite de droite}

\begin{defi}{(rappel) Equation réduite de droite}{}
   Soit deux coefficients $m, p\in\R$. On appelle \textbf{équation reduite de droite}, l'équation suivante d'inconnue $(x;y)$
   $$y=mx+p$$

   Le coefficient $m$ est appelé \textbf{coefficient directeur}. Le coefficient $p$ est appelé \textbf{ordonnée à l'origine}.
\end{defi}

\begin{prop}{Equation réduite}{}
   Soit la droite $\mathcal{D}$ d'équation cartésienne $ax+by+c=0$, avec $a,b,c\in\R$ et $(a;b)\ne(0;0)$.
   \medskip
   La droite $\mathcal{D}$ est non parallèle à l'axe des ordonnées (non verticale), si et seulement s'il existe des coefficients $m,p\in\R$ tels que la droite $\mathcal{D}$ admette $y =mx+p$ comme équation réduite.
   \medskip
   De plus, $m=\frac{-a}{b}$ et $p=\frac{-c}{b}$.
\end{prop}

\begin{dem}{}{}
   \begin{itemize}
      \item Commençons par le sens direct ($\implies$). On suppose que la droite $\mathcal{D}$ n'est pas verticale. 
      
      D'après la propriété $\ref{propo:droite_verticale}$, puisque la droite $\mathcal{D}$ n'est pas verticale, alors $b\troum{\ne 0}$
      \begin{reponse}
         Donc on peut réécrire l'équation sous la forme $\frac{a}{b}x+y+\frac{c}{a}=0$.

         On encore, $y=mx+p$ avec $m=\frac{-a}{b}$ et $p=\frac{-c}{b}$
      \end{reponse}
      \item Prouvons le sens réciproque ($\impliedby$). On suppose qu'il existe des coefficients $m,p\in\R$ tels que la droite $\mathcal{D}$ admette $y =mx+p$ comme équation réduite.
      \begin{reponse}
         
         On veut prouver que $\mathcal{D}$ n'est pas verticale. Il suffit de trouver deux points de $\mathcal{D}$ d'abscisses différentes.
         
         Or $(0;p)\in\mathcal{D}$ et $(1;m+p)\in\mathcal{D}$. Donc la droite n'est pas verticale.
      \end{reponse}
   \end{itemize}
   
\end{dem}

\begin{prop}{(rappel) Calcul du coefficient directeur}{}
   Soient $A(x_A;y_A)$ et $B(x_B;y_B)$ deux points d'une droite $\mathcal{D}$, le coefficient directeur se calcule grâce à la formule : $$m=\frac{y_B-y_A}{x_B-x_A}=\frac{\Delta y}{\Delta x}=\frac{\text{variation ordonnée}}{\text{variation abscisse}}$$
\end{prop}

\begin{methode}{Déterminer l'équation réduite...}{}
\textbf{A partir des points $A\coordpp{x_A}{y_A}$ et $B\coordpp{x_B}{y_B}$}

\medskip
Nous prendrons l'exemple de $A(2;-1)$ et $B(-1;5)$.
La droite $D$ n'est pas parallèle à l'axe des ordonnées, elle a donc pour équation $y=mx+p$.

\medskip
\hspace{5mm}$m=\troum{\dfrac{y_B-y_A}{x_B-x_A}=\dfrac{5+1}{-1-2}=-2}$

\medskip
\noindent $\mathcal{D}$ passe par le point $A(2;-1)$ donc les coordonnées du point $A$ vérifie l'équation de la droite $\mathcal{D}$ d'où :

\begin{reponse}
   \hspace{5mm}\phantom{$\Longleftrightarrow$} $y_A=mx_A+p$
   
   \hspace{5mm}$\Longleftrightarrow y_A=-2x_A+p$
   
   \hspace{5mm}$\Longleftrightarrow -1=-2\times2+p$
   
   \hspace{5mm}$\Longleftrightarrow p=3$
\end{reponse}
   
\medskip
La droite $\mathcal{D}$ a pour équation réduite : $y=-2x+3$
\end{methode}


\begin{prop}{Vecteur directeur et coefficient directeur}{}
   Soit deux réel $m,p\in\R$. Soit une droite non verticale $\mathcal{D}$ du plan $\R^2$ d'équation réduite $y = mx+p$.
   
   \begin{itemize}
      \item 
   Le vecteur $\vecu \coordvp{1}{m}$ est un vecteur directeur de la droite $\mathcal{D}$.
      \item Réciproquement, si $\vecu\coordvp{x}{y}$ est un vecteur directeur, alors $m=\frac{y}{x}$
   \end{itemize}
\end{prop}

% \begin{dem}{}{}
%    La démonstration est similaire à la propriété \ref{propo:vecteur_directeur_eq_cartesienne}
% \end{dem}
%\begin{ex}{}{}
%Pour chacune des équations suivantes, déterminons, s'il existe, le coefficient directeur et l'ordonnée à l'origine de la droite :
   
%\renewcommand{\arraystretch}{1.5}
%\begin{tabular}{p{0.2cm}p{4cm}p{3cm}p{3cm}p{3cm}}
%&  & Équation réduite & Coef. directeur & Ordonnée à l'origine \\
%& $5x+y=-2$ & $y=-5x-2$ & $m=-5$ & $p=-2$ \\
%& $2x-3y=1$ & $y=\frac{2}{3}x-\frac{1}{3}$ & $m=\frac{2}{3}$ & $p=-\frac{1}{3}$ \\
%& $2x+3y=7-5x+3y$ & $x=1$ & \multicolumn{2}{l}{Droite parallèle à l'axe des ordonnées}
%\end{tabular}
%\end{ex}


\vspace{-5mm}
\section{Droites parallèles}

\begin{prop}{}{}
   Les trois propositions suivantes sont équivalentes :

   \begin{enumerate}[label=(\roman*)]
      \item 
      Les droites $\mathcal{D}$ et $\mathcal{D'}$ sont parallèles.
      \item 
      Les droites $\mathcal{D}$ et $\mathcal{D'}$ ont le même coefficient directeur (ou sont toutes deux verticales).
      \item 
      Les droites $\mathcal{D}$ et $\mathcal{D'}$ admettent un même vecteur directeur (leurs vecteurs directeurs sont colinéaires).
   \end{enumerate}
\end{prop}

\newpage

% \begin{ex}{}{}
% Déterminer une équation de la droite $\mathcal{D}$ parallèle à la droite $\mathcal{D'}$ d'équation $y=2x-3$ passant par le point $A(1;5)$.

% \medskip
% \noindent La droite $D$ n'est pas parallèle à l'axe des ordonnées, elle a donc pour équation $y=mx+p$.

% \medskip
% \noindent $\mathcal{D}$ est parallèle à $\mathcal{D}'$ donc $m_{\mathcal{D}}=m_{\mathcal{D}'}=2$

% \medskip
% \noindent $\mathcal{D}$ passe par le point $A(1;5)$ d'où :

% \medskip
% \hspace*{0.2cm} $y_A=mx_A+p$
              
% \hspace*{0.2cm} $5=2\times 1+p$
              
% \hspace*{0.2cm} $p=3$

% \medskip
% \noindent La droite $\mathcal{D}$ a pour équation réduite $y=2x+3$ 
% \end{ex}

\vspace{-5mm}
\section{Représentation graphique}

Dans cette section, on se place dans le plan $\R^2$ muni d'un repère $\oij$.
Comment tracer la représentation graphique d'une droite ?

\noindent \begin{minipage}{8cm}
\begin{methode}{Tracer une droite...}{}
   \textbf{A partir de l'équation cartésienne ou réduite}
\medskip
\begin{itemize}
\item Choisir deux abscisses $x_{1}$ et $x_{2}$,\medskip
\item Calculer $y_{1}$ et $y_{2}$ avec l'équation,\medskip
\item Placer les points $A(x_1;y_1 )$ et $B(x_2;y_2)$ \medskip
\item Tracer la droite $(AB)$
\end{itemize}
\end{methode}
\end{minipage}
\hspace{5mm}
\noindent \begin{minipage}{8cm}
   \begin{methode}{Tracer une droite...}{}
      \textbf{A partir d'un point $A$ et du vecteur directeur directeur $\vecu$}
      \medskip
      \begin{itemize}
      \item Placer le point $A$,\medskip
      \item A partir de ce point, tracer le vecteur directeur $\vecu$ ce qui nous donne un second point $M$\medskip
      \item Tracer la droite $(MP)$.
      \end{itemize}
      \end{methode}
\end{minipage}

\begin{rem}{}{}
   Dans le cas, où l'on connait l'équation cartésienne ou réduite, on est bien sûr capable de trouver des vecteurs directeurs (respectivement $\vecu\coordvp{-b}{a}$ et $\vecu\coordvp{1}{m}$), et on peut utiliser la deuxième méthode.
\end{rem}

\begin{ex}{}{}
   Représentons graphiquement la droite d'équation $y=\dfrac{2}{3}x-2$ avec les différentes méthodes.
\begin{tikzpicture}[line cap=round,line join=round,>=triangle 45,x=0.65cm,y=0.65cm]
   \begin{axis}[
         x=1.0cm,y=1.0cm,
         axis lines=middle,ymajorgrids=true,
         xmajorgrids=true,
         xmin=-5.5,xmax=5.5,
         ymin=-5.5,ymax=5.5,
         xtick={-5,-4,...,5},ytick={-5,-4,...,5},
         ]
      \end{axis}
   \end{tikzpicture}
\end{ex}

% \begin{minipage}{10.5cm}
% Représentons graphiquement la droite d'équation $y=\dfrac{2}{3}x-2$ avec la méthode des coordonnées de points.

% \begin{itemize}
% \item Choix de $x_{1}$ et $x_{2}$. Le mieux est de prendre deux valeurs pouvant se placer dans le repère tracé (ici entre $-4$ et $4$) et relativement écartées afin d'avoir une meilleur précision de tracé.

% Je choisis ici $x_1=-3$ et $x_2=3$ afin de simplifier les calcul de $y_{1}$ et $y_{2}$
% \item $y_{1} = \dfrac{2}{3}x_{1}-2$ et $y_{2} = \dfrac{2}{3}x_{2}-2$

% Donc : $y_{1} = \dfrac{2}{3}\times (-3)-2$ et $y_{2} = \dfrac{2}{3}\times (3)-2$

% Donc : $y_{1} = -4$ et $y_{2} = 0$
% \item On place les points $A(x_1;y_1 )$ et $B(x_2;y_2)$ (en bleu sur le graphique)
% \item On la droite $(AB)$ (en rouge sur le graphique)
% \end{itemize}
% \end{minipage}
% \hspace{5mm}
% \begin{minipage}{5.5cm}
% \begin{center}
% \begin{tikzpicture}[line cap=round,line join=round,>=triangle 45,x=0.7cm,y=0.7cm,scale=0.8]
% \draw [color=lightgray,, xstep=0.7cm,ystep=0.7cm] (-4.5,-4.5) grid (4.5,4.5);
% \draw[->,color=black] (-4.5,0.) -- (4.5,0.);
% \foreach \x in {-4.,-3.,-2.,-1.,1.,2.,3.,4.}
% \draw[shift={(\x,0)},color=black] (0pt,2pt) -- (0pt,-2pt) node[below] {\footnotesize $\x$};
% \draw[->,color=black] (0.,-4.5) -- (0.,4.5);
% \foreach \y in {-4.,-3.,-2.,-1.,1.,2.,3.,4.}
% \draw[shift={(0,\y)},color=black] (2pt,0pt) -- (-2pt,0pt) node[left] {\footnotesize $\y$};
% \draw[color=black] (0pt,-10pt) node[right] {\footnotesize $0$};
% \clip(-4.5,-4.5) rectangle (4.5,4.5);
% \draw[line width=1.5pt,color=red,smooth,samples=100,domain=-4.5:4.5] plot(\x,{0-2.0+0.66666666667*(\x)});
% \draw (-3,-4) node{\color{blue} \Large{$\times$}};
% \draw (3,0) node{\color{blue} \Large{$\times$}};
% \end{tikzpicture}
% \end{center}
% \end{minipage}
% \end{ex}

% \vspace{-5mm}
% \section{Résolution graphique des systèmes d'équation}


\newpage
\section*{Exercices}
\begin{exercices}
	\sectionexo{Equation cartésienne}
	\begin{bookexo}{34 p.136}
      Déterminer également l'ordonnée du point d'abscisse $x=3$.
   \end{bookexo}
	\sectionexo{Vecteur directeur}
   \begin{bookexo}{37 p.136}\end{bookexo}
   \begin{bookexo}{44 p.137}\end{bookexo}
   \begin{bookexo}{61 p.139}\end{bookexo}
	\sectionexo{Equation réduite}
   \begin{bookexo}{40 p.136}\end{bookexo}
	\sectionexo{Représentation graphique}
   \begin{bookexo}{50 p.138}\end{bookexo}
   \begin{bookexo}{54 p.138}\end{bookexo}
   \begin{bookexo}{58 p.138}\end{bookexo}
	\sectionexo{Systèmes d'équation}
   \begin{bookexo}{73 p.165}\end{bookexo}

\end{exercices}
\end{document}

